\section{Lift nonadic of the unit of Pell}

For \(k\ge1\), let the unramified quadratic ring be
\begin{equation}
\mathcal R_k=(\Z/3^k\Z)[r]/(r^2-2),
\label{eq:pell-ring}
\end{equation}
and let

\[
u=3+2r.
\]

The conjugation \(r\mapsto-r\) gives
\[
N_{\mathcal R_k/(\Z/3^k\Z)}(u)
=(3+2r)(3-2r)=9-8=1,
\]
so that \(u\) is a unit of norm one. Moreover,
\[
u^2=17+12r,\qquad
u^4-1=576+408r=3(192+136r).
\]
The second factor is a unit modulo three because it reduces to \(r\ne0\). Hence
the valuation \(3\)-adic satisfies
\(v_3(u^4-1)=1\). The lifting lemma for a unit congruent to one
modulo three yields, for every \(j\ge0\),
\[
v_3\!\left(u^{4\cdot3^j}-1\right)=j+1.
\]
Since \(u\bmod3=2r\) and \((2r)^2=-1\), its order modulo three is four.
Consequently,

\begin{equation}
\operatorname{ord}(u)=4\cdot3^{k-1},
\label{eq:pell-order}
\end{equation}

and the profinite closure of the compatible cyclic groups is
\begin{equation}
\varprojlim_k\langle u\bmod3^k\rangle
\simeq C_4\times\mathbb Z_3.
\label{eq:pell-profinite}
\end{equation}
Its Archimedean realization is

\[
3+2\sqrt2=(1+\sqrt2)^2.
\]

This construction proves the conservation of nonadic depth
by a hyperbolic unit. The units \(3+2\sqrt2\),
\(2+\sqrt3\), and \(30+\sqrt{899}\) form a Pell family with
distinct provenances; any identification among them requires
explicit maps between the corresponding extensions.

